The clearest forecasting result is the value of allowing industries to differ. Separate industry regressions reduce average error relative to the original pooled specification and improve the previously weak 50--99 employee band. Regularized pooled models also improve performance. The unrestricted full dictionary performs substantially worse, and removing quantitative curvature from ridge has little effect. These findings favor targeted heterogeneity, historical smoothing, and coefficient control over indiscriminate expansion of the regression equation.

Nonlinear subset selection identifies industry-dependent slopes and a curved association between expectations and future published use. The selected expectation form changes across validation dates and does not deliver an aggregate improvement over the baseline. The final subset gives a decreasing conditional relationship with expected use over the plotted ranges, further limiting a structural interpretation. The statistical conclusion is therefore limited: the candidate dictionary contains nonlinear descriptions, but their incremental predictive value has not been established. A subsequent preregistered comparison on genuinely future releases would provide stronger evidence.

For industry comparisons, finance, Information, and professional services form a relatively stable leading growth group. Their exact ordering is uncertain. Outreach organizations could use these results to distinguish broad adoption trajectories while avoiding precise league tables. Higher adoption or faster growth does not establish unmet need, accessibility of assistance, or the benefit of an intervention. Resource-allocation decisions require additional information on barriers, costs, and program outcomes.

The interpretation of uncertainty remains conditional. Estimates arise from a short, noisy panel with recurring groups, shared date shocks, and a rotating business sample. The six-release block bootstrap has few effective temporal blocks. It is not a substitute for Census replicate weights or full design-based covariance information. Standard errors address sampling variability, while nonresponse and interpretation changes can introduce additional error. Screening excludes groups with frequent missing estimates, and the results may not transfer to those groups. Equal weighting across retained cells also differs from weighting industries by their business populations.

The event study supplies a useful chronology but cannot attribute adoption changes to specific releases. Announcements are common exposures, cluster early in the sample, and overlap with a strong adoption trend. Early releases lack comparable pre-event data. A stronger design would link firm or industry exposure to particular tools, specify comparison groups before inspecting outcomes, and assess pre-event trends over a longer consistent series.

Finally, the capacity experiment shows why parameter count alone is an inadequate guide to forecasting. Minimum-norm interpolation produces extreme errors in this noisy setting, although error declines after the interpolation peak. Ridge stabilizes a much larger feature representation through shrinkage. These results are specific to the tested random-feature distribution, fixed bandwidth, penalty, origins, and seeds; they do not establish a universal double-descent curve for BTOS or for machine learning generally.
